\section{问题一：预热平衡阶段的温湿耦合场}

\subsection{模型的建立}

中药材为半径 $R=2.0$~cm、长 $25$~cm 的圆柱。烘房热风沿轴向与周向均匀作用，
而药材长径比超过 $12$，两端面的轴向输运相对于径向可以忽略，
因此我们取\textbf{一维径向轴对称}模型，即第五章式~\eqref{eq:energy} 与式~\eqref{eq:mass}：

\begin{equation}
\begin{cases}
\rho c_p\dfrac{\partial T}{\partial t}=\dfrac{1}{r}\dfrac{\partial}{\partial r}
\left(k\,r\,\dfrac{\partial T}{\partial r}\right)\\[16pt]
\dfrac{\partial C}{\partial t}=\dfrac{1}{r}\dfrac{\partial}{\partial r}
\left(D\,r\,\dfrac{\partial C}{\partial r}\right)
\end{cases}
\label{eq:q1gov}
\end{equation}

边界取式~\eqref{eq:bc} 的第三类对流条件，在 $r=R$ 处有

\[
\begin{cases}
-k\,\dfrac{\partial T}{\partial r}=h\left(T-T_\infty\right)\\[16pt]
-D\,\dfrac{\partial C}{\partial r}=h_m\left(C-C_\infty\right)
\end{cases}
\]

在 $r=0$ 处由轴对称性给出零梯度。初值取

\[
\begin{cases}
T=28~^\circ\text{C}\\[16pt]
C=2.55~\text{kg/kg}
\end{cases}
\]

且全域均匀。环境温湿度 $T_\infty(t)$、$C_\infty(t)$ 取附件 1 的实测序列并线性插值——
题目要求的 $t=100$~s 不是附件 1 的采样点，必须插值才能得到。
附录 2 给出的本问物性汇总于表~\ref{tab:q1props}。

\begin{table}[H]
  \centering
  \caption{问题一所用物性参数}\label{tab:q1props}
  \small
  \begin{tabular}{@{}llll@{}}
    \toprule
    物性 & 符号 & 取值 & 单位 \\
    \midrule
    药材密度 & $\rho$ & $820$ & kg/m$^3$ \\
    比热容 & $c_p$ & $2600$ & J/(kg$\cdot$K) \\
    导热系数 & $k$ & $0.36$ & W/(m$\cdot$K) \\
    对流传热系数 & $h$ & $25$ & W/(m$^2\cdot$K) \\
    对流传质系数 & $h_m$ & $8\times10^{-7}$ & m/s \\
    水分扩散系数 & $D$ & $7\times10^{-9}e^{-0.89/C}$ & m$^2$/s \\
    \bottomrule
  \end{tabular}
\end{table}

质量方程按题目明文采用 Fick 标准形式，$\rho$、$c_p$、$k$ 只进入能量方程。
注意到本问物性为常数，温度不进入扩散系数、含水率也不进入热物性，
因此式~\eqref{eq:q1gov} 的两个方程彼此独立，可以分别求解。

\subsection{模型的求解}

本问沿用第五章的控制方程与顶点有限体积离散（$N=800$、$\Delta r=0.025$~mm，
见式~\eqref{eq:semi} 与式~\eqref{eq:surf}），本节的增量只在时间推进。

$\theta$ 加权格式对本问同样适用，但本问物性为常数、两个方程都是常系数线性方程，
Jacobian 为常矩阵，故 BDF 每步只解一次线性方程组、分解可跨步复用，自适应步长几乎无额外代价。
因此主方法采用自适应隐式多步 BDF（$\mathrm{rtol}=10^{-9}$、$\mathrm{atol}=10^{-11}$）：
突变段自动降为一阶小步、平滑段升为高阶大步，既相当于自动完成后向 Euler 启动，
又使时间误差由容差直接量化，本问的时间精度证据正是容差加密。

按第五章的统一验证设计，本问另用一套\textbf{格心有限体积}实现独立复算（表面值由 Robin 关系回算、
时间改用固定步长 Crank--Nicolson），与主方法在表面值来源与时间推进上均不相同，
属于独立核算而非同一算法的重复实现；逐点对照见本章数值检验一节，求解流程见图~\ref{fig:flowq1}。

\begin{figure}[H]
  \centering
  \includegraphics[width=0.90\textwidth]{fig08_flow_q1.pdf}
  \caption{问题一的求解流程}
  \label{fig:flowq1}
\end{figure}

\subsection{结果与分析}

\input{paper/tables/table1_q1_temperature.tex}
\begin{table}[H]
  \centering
  \caption{预热平衡阶段的干基水分浓度（kg/kg）}
  \label{tab:table2_q1_moisture}
  \small
  \begin{tabular}{lccccc}
    \toprule
    时间/s & $r=0$ & $0.5$~cm & $1.0$~cm & $1.5$~cm & $2.0$~cm \\
    \midrule
    100 & 2.5500 & 2.5500 & 2.5500 & 2.5500 & 2.2471 \\
    300 & 2.5500 & 2.5500 & 2.5500 & 2.5492 & 2.0508 \\
    600 & 2.5500 & 2.5500 & 2.5500 & 2.5353 & 1.8770 \\
    900 & 2.5500 & 2.5500 & 2.5497 & 2.5045 & 1.7547 \\
    1200 & 2.5500 & 2.5500 & 2.5482 & 2.4646 & 1.6586 \\
    1500 & 2.5500 & 2.5499 & 2.5445 & 2.4206 & 1.5787 \\
    1800 & 2.5500 & 2.5497 & 2.5383 & 2.3755 & 1.5102 \\
    \bottomrule
  \end{tabular}
\end{table}


\textbf{在 $1800$~s 时刻，中心温度为 \qoneTcenteroneeightzerozeros，表面温度为 \qoneTsurfaceoneeightzerozeros；
中心干基含水率仍为 \qoneCcenteroneeightzerozeros，而表面已降至 \qoneCsurfaceoneeightzerozeros。}
表~\ref{tab:table1_q1_temperature} 与表~\ref{tab:table2_q1_moisture} 即题目要求的两张表，
其最后一行给出这一时刻的读数：温度自中心向表面递增、含水率自中心向表面递减，中心为全域最大值。
图~\ref{fig:q1} 进一步给出七个报告时刻的剖面演化与中心、表面时程：温度剖面随报告时刻整体上移、
在靠近表面处才变陡，含水率的下凹则始终局限于表层，中心含水率在整个窗口内几乎不变。
图~\ref{fig:q1map} 以附录图的形式给出预热阶段温湿场的完整时空分布。



\begin{figure}[H]
  \centering
  \includegraphics[width=0.98\textwidth]{fig02_preheat_fields.pdf}
  \caption{预热阶段的温度与水分场}
  \label{fig:q1}
\end{figure}

\begin{figure}[H]
  \centering
  \includegraphics[width=0.95\textwidth]{figA1_q1_field_map.pdf}
  \caption{预热阶段温湿场（附录图）}
  \label{fig:q1map}
\end{figure}

上述结果呈现出本阶段最基本的规律：\textbf{升温已经接近完成，而失水仍局限于很薄的一层}。
从表~\ref{tab:table2_q1_moisture} 可见，含水率在距表面 $0.5$~cm 以内的范围内已明显下降，
而在中心附近几乎没有变化。这一现象的原因在第五章已由特征时间给出：
水分扩散的特征时间 $R^2/D\approx22$~h，远长于本阶段的 $30$~min，
即本阶段只走完了全部扩散过程的约 $2\%$，水分的重新分布尚不足以影响内部。
可见在这个时间尺度上，内部含水率约为初值是合理的，这也说明本文建立的模型与
该阶段的物理过程是相符的。



\subsection{数值检验与模型适用性}



结果的可信度由四类相互独立的证据确认，全部检验量汇总于表~\ref{tab:verify}。

\begin{enumerate}[label=(\arabic*), leftmargin=2.4em, itemsep=2pt, topsep=3pt, parsep=0pt]
  \item \textbf{与解析解对照}：在常物性、固定半径的极限下，本文实现与圆柱扩散问题的
        Bessel 级数解析解相比，最大绝对误差仅为 \qoneanalyticalmaxabserr，说明离散格式本身是正确的。
  \item \textbf{两项加密实验}：网格由 $400$ 加密到 $800$，报告点上的最大差为 \qonegridrefinementdiff；
        求解器容差由 $10^{-8}$ 收紧到 $10^{-9}$，表值变化为 \qonetimeaccuracytolerancediff。
        两者都已低于题目要求的四位小数，可见所报数值已达到网格与步长无关的水平。
  \item \textbf{守恒性}：把离散质量守恒写成式~\eqref{eq:conservation} 那样的恒等式后逐时步核验，
        本问的相对残差为 \qonemassbalancerel，达到机器精度量级。
  \item \textbf{独立复算}：本文独立编写的格心实现与顶点格式在全部 $35$ 个报告点上的
        收敛极限相差 \qoneindependentimplementationsdiff，两法逐点一致到四位小数的舍入边界。
\end{enumerate}

\textbf{四类证据一致证明本问报告数值的可信度。}
